\chapter{Trabajo futuro}

Aunque el proyecto cumple satisfactoriamente con los objetivos propuestos, existen diversas mejoras que podrían incorporarse para incrementar el nivel de realismo y ampliar las funcionalidades de la simulación.

\begin{itemize}

\item Implementar sombras dinámicas generadas por los relámpagos utilizando mapas de sombras (Shadow Mapping).

\item Incorporar efectos de niebla y volumetría para representar condiciones climáticas más realistas durante tormentas intensas.

\item Agregar animaciones de balanceo en los árboles en función de la intensidad y dirección del viento.

\item Implementar múltiples tipos de precipitación, como nieve, granizo o lluvia ligera, reutilizando la arquitectura del sistema de partículas.

\item Incorporar sonidos ambientales adicionales, incluyendo viento, truenos con diferentes intensidades y sonidos producidos por la lluvia al impactar distintas superficies.

\item Desarrollar un sistema de clima dinámico capaz de alternar automáticamente entre cielo despejado, lluvia, tormenta y otros estados meteorológicos.

\item Optimizar el renderizado utilizando técnicas modernas de OpenGL como Instanced Rendering para representar un número considerablemente mayor de partículas y vegetación.

\item Implementar un sistema de terreno basado en alturas (Height Maps) que permita representar montañas, colinas y desniveles en lugar de un plano completamente uniforme.

\item Incorporar modelos tridimensionales para algunos elementos del escenario, combinándolos con billboards para mantener un adecuado equilibrio entre calidad visual y rendimiento.

\item Desarrollar un sistema de iluminación física (PBR) que mejore el comportamiento de los materiales frente a la luz generada por los relámpagos.

\item Añadir soporte para cámaras completamente libres controladas mediante el mouse, permitiendo explorar la escena desde cualquier ángulo.

\item Incorporar un ciclo completo de día y noche con variaciones automáticas de iluminación y condiciones climáticas.

\end{itemize}

La arquitectura modular implementada durante este proyecto facilita la incorporación de estas mejoras, ya que cada componente se encuentra desacoplado del resto del sistema, permitiendo extender la aplicación sin necesidad de modificar significativamente la estructura existente.